\scanpage{263}
\mspage{133 quater}

et cette fois c'est $\mathscr H'$ qui est solution du pb.
universel ($\mathrm{Cat}\unsure{-|\to}\mathscr H'$ transformant flèches de
$W$ en iso). Donc dans ce cas $\mathscr H'\unsure{-|\to}\mathrm{Hot}_{W}$ est
un iso, et $\mathrm{Hot}_{W}$ est équiv. à $e$. Dans les
deux cas, $\mathrm{Hot}_{W}^{*}$ est équivalente à $e$.

Donc, la prop. 9.12. résulte du cor. \unsure{9,13.|9.13.}, qu'il
s'agit à présent de démontrer.

Soit $f:\deleted{\unsure{X}}I\to\deleted{\unsure{Y}}J$ \inserted{de $\mathrm{Cat}$} telle que $\pi_0(f)$ ne soit pas injective.
On a vu dans la dém. de 9.7 comment on en
déduit (en choisissant deux éléments $i_0,i_1$ dans deux
comp. connexes distinctes de $I$, telles que
\deleted{\unsure{$j_1$}} $j_0=f(i_0)$ et $j_1=f(i_1)$ \unsure{'} soient dans une même
composante connexe de $\unsure{J|D}$) comment on déduit
que pour $X,Y$ dans $\mathrm{Cat}$, deux flèches $f_0,f_1:X\mathrel{\unsure{=|\rightrightarrows}}Y$
sont toujours égales dans \deleted{\unsure{$W$}} $\mathrm{Hot}_{W}$. (cf lemme 9.10.)

Soit maintenant

\[
 f:I\unsure{-|\to}J
\]

dans $\mathrm{Cat}$ telle que $\pi_0(f)$ ne soit pas surjective,
donc $J=J_0\amalg J_{\unsure{1|0}}$, avec $f$ se factorisant par
$J_0$, et $J_1\ne\varnothing$. Considérons alors la cat. discrète
à deux objets,

\[
 \unsure{\varepsilon|E}=\{0,1\},
\]

et deux \deleted{\unsure{appl}} flèches

\[
 g_0,g_1:J\rightrightarrows\varepsilon,
\]

$g_0$ étant \unsure{constante} de valeur $0$, et $g_1$ étant
égale à $0$ sur $J_0$, à $1$ sur $J_1$. On a

\[
 g_0f=g_1f\quad\text{\unsure{et|d}}
\]

donc, comme $f\in W$,

\[
 g_0=g_1\qquad\text{dans }\mathrm{Hot}_{W}.
\]

On en conclut, en prenant un point $j_1$ dans $J_1$,
et l'inclusion $\alpha_{\unsure{j_1|1}}:e\longrightarrow J$ correspondante, que
% archive sheet number: 262
